\documentclass[../../main.tex]{subfiles}% 注意这里的文件路径不能用 ./main.tex ,否则用latexmk编译子文件会报错
\graphicspath{{\subfix{./image/}}{\subfix{../../image/}}} % 指定图片引用路径,后续可以直接使用图片文件名
% 如果图片存放在上级目录,则使用 ../../image/ 作为备用路径

% 例如:
% \begin{figure}[H]
% \centering
% \includegraphics[width=0.3\textwidth]{图.png}
% \caption{}
% \label{figure:图}
% \end{figure}
% 注意:上述\label{}一定要放在\caption{}之后,否则引用图片序号会只会显示??.

\begin{document}

\section{习题}

\begin{example}
按定义求下列函数的 Fourier 变换：
\begin{enumerate}[(1)]
\item $f(x)=\begin{cases}
0, & |x|>a,\\
|x|, & |x|\leqslant a
\end{cases}\quad (a>0)$；
\item $f(x)=\begin{cases}
0, & |x|>a,\\
1-\frac{|x|}{a}, & |x|\leqslant a
\end{cases}\quad (a>0)$；
\item $f(x)=\begin{cases}
0, & |x|>a,\\
\sin\lambda_0 x, & |x|\leqslant a
\end{cases}\quad (a>0,\ \lambda_0>0)$；
\item $f(x)=e^{-a|x|}\quad (a>0)$；
\item $f(x)=\cos x\,e^{-a|x|}\quad (a>0)$。
\end{enumerate}
\end{example}

\begin{example}
利用 Fourier 变换的性质求下列函数的 Fourier 变换：
\begin{enumerate}[(1)]
\item $f(x)=e^{-a^2x^2}\cos bx$；
\item $f(x)=e^{-a^2x^2}\sin bx$；
\item $f(x)=xe^{-a^2x^2}$；
\item $f(x)=x^2e^{-a^2x^2}$；
\item $f(x)=e^{-a|x|}\cos bx$；
\item $f(x)=e^{-a|x|}\sin bx$；
\item $f(x)=xe^{-a|x|}$；
\item $f(x)=x^2e^{-a|x|}$；
\item $f(x)=\begin{cases}
e^{-ax}, & x>0,\\
0, & x<0
\end{cases}$；
\item $f(x)=\begin{cases}
xe^{-ax}, & x>0,\\
0, & x<0
\end{cases}$。
\end{enumerate}
\end{example}

\begin{example}
求下列函数的 Fourier 逆变换：
\begin{enumerate}[(1)]
\item $F(\lambda)=e^{-a^2\lambda^2 t}$，其中 $t>0$ 为参数，$a>0$ 为常数；
\item $F(\lambda)=e^{(-a^2\lambda^2+\mathrm{i}b\lambda+c)t}$，其中 $t>0$ 为参数，$a\in\mathbb{R}_+$，$b,c\in\mathbb{R}$ 为常数；
\item $F(\lambda)=e^{-|\lambda|y}$，其中 $y>0$ 为参数。
\end{enumerate}
\end{example}

\begin{example}
利用 Fourier 变换求解下列定解问题：
\begin{enumerate}[(1)]
\item $\begin{cases}
u_t-a^2u_{xx}=0, & (x,t)\in\mathbb{R}\times\mathbb{R}_+,\\
u(x,0)=\varphi(x), & x\in\mathbb{R},
\end{cases}$
其中 $a>0$ 为常数，$\varphi\in C(\mathbb{R})\cap L^\infty(\mathbb{R})$；
\item $\begin{cases}
u_t-a^2u_{xx}+bu_x+cu=0, & (x,t)\in\mathbb{R}\times\mathbb{R}_+,\\
u(x,0)=\varphi(x), & x\in\mathbb{R},
\end{cases}$
其中 $a>0$ 为常数，$b,c\in\mathbb{R}$ 是常数，$\varphi\in C(\mathbb{R})\cap L^\infty(\mathbb{R})$；
\item $\begin{cases}
u_t-a^2u_{xx}=f(x,t), & (x,t)\in\mathbb{R}\times\mathbb{R}_+,\\
u(x,0)=0, & x\in\mathbb{R},
\end{cases}$
其中 $a>0$ 为常数，$f$ 为已知函数。
\end{enumerate}
\end{example}
\begin{solution}
对 $x$ 作 Fourier 变换，记
\begin{align*}
\hat u(\xi,t)=\int_{-\infty}^{+\infty}u(x,t)e^{-\mathrm{i}\xi  x}\,\mathrm{d}x.
\end{align*}
\begin{enumerate}[(1)]
\item 对原方程 $u_t-a^2u_{xx}=0$ 两边作Fourier变换，利用\hyperref[proposition:Fourier 变换的基本性质1-2]{Fourier变换的微商性质}得
\begin{align*}
\hat{u}_t-a^2\hat{u}_{xx}=\hat{u}_t-a^2(\mathrm{i}\xi )^2\hat{u}=\hat{u}_t+a^2\xi ^2\hat{u}=0.
\end{align*}
这是关于 $t$ 的一阶常微分方程，解得
\begin{align*}
\hat u(\xi,t)=\hat u(\xi,0)e^{-a^2\xi^2t}.
\end{align*}
由初值条件 $u(x,0)=\varphi(x)$，有 $\hat u(\xi,0)=\hat\varphi(\xi)$. 故
\begin{align*}
\hat u(\xi,t)=\hat\varphi(\xi)e^{-a^2\xi^2t}.
\end{align*}
作逆 Fourier 变换，得
\begin{align*}
u(x,t)=\frac{1}{\sqrt{2\pi}}\int_{-\infty}^{+\infty}\hat\varphi(\xi)e^{-a^2\xi^2t}e^{\mathrm{i}\xi  x}\,\mathrm{d}\xi.
\end{align*}

\item 对方程 $u_t-a^2u_{xx}+bu_x+cu=0$ 两边作Fourier变换，利用\hyperref[proposition:Fourier 变换的基本性质1-2]{Fourier变换的微商性质}得
\begin{align*}
\hat u_t+a^2\xi^2\hat u+\mathrm{i}b\xi\hat u+c\hat u=0,
\end{align*}
即
\begin{align*}
\hat u_t+(a^2\xi^2+\mathrm{i}b\xi+c)\hat u=0.
\end{align*}
解得
\begin{align*}
\hat u(\xi,t)=\hat u(\xi,0)\exp\left[-(a^2\xi^2+\mathrm{i}b\xi+c)t\right].
\end{align*}
由初值条件，$\hat u(\xi,0)=\hat\varphi(\xi)$，故
\begin{align*}
\hat u(\xi,t)=\hat\varphi(\xi)e^{-ct}e^{-a^2\xi^2t}e^{-\mathrm{i}b\xi t}.
\end{align*}
作逆 Fourier 变换得
\begin{align*}
u(x,t)=\frac{1}{\sqrt{2\pi}}\int_{-\infty}^{+\infty}\hat\varphi(\xi)e^{\left[-(a^2\xi^2+\mathrm{i}b\xi+c)t\right]}e^{\mathrm{i}\xi x}\,d\xi.
\end{align*}

\item 对原方程$u_t-a^2u_{xx}=f(x,t)$两边作 Fourier 变换，利用\hyperref[proposition:Fourier 变换的基本性质1-2]{Fourier变换的微商性质}得
\begin{align*}
\hat u_t+a^2\xi^2\hat u=\hat f(\xi,t),
\end{align*}
其中 $\hat f(\xi,t)$ 是 $f(x,t)$ 关于 $x$ 的 Fourier 变换。初值条件 $u(x,0)=0$ 给出 $\hat u(\xi,0)=0$. 解此一阶线性常微分方程得
\begin{align*}
\hat u(\xi,t)=\int_0^t \hat f(\xi,\tau)e^{-a^2\xi^2(t-\tau)}\,\mathrm{d}\tau.
\end{align*}
作逆 Fourier 变换得
\begin{align*}
u(x,t)=\frac{1}{\sqrt{2\pi}}\int_{-\infty}^{+\infty}\left[\int_0^t \hat f(\xi,\tau)e^{-a^2\xi^2(t-\tau)}\,d\tau\right]e^{\mathrm{i}\xi x}\,d\xi.
\end{align*}
\end{enumerate}
\end{solution}

\begin{example}
设 $\Phi(\lambda)=\hat{\varphi}(\lambda)$，求下列函数的 Fourier 逆变换：
\begin{enumerate}[(1)]
\item $F(\lambda)=\Phi(\lambda)\cos a\lambda t$，其中 $t>0$ 为参数，$a>0$ 为常数；
\item $F(\lambda)=\Phi(\lambda)\frac{\sin a\lambda t}{a\lambda}$，其中 $t>0$ 为参数，$a>0$ 为常数。
\end{enumerate}
\end{example}

\begin{example}
利用 Fourier 变换求解波动方程初值问题
\begin{align*}
\begin{cases}
u_{tt}-a^2u_{xx}=0, & (x,t)\in\mathbb{R}\times\mathbb{R}_+,\\
u(x,0)=\varphi(x), & x\in\mathbb{R},\\
u_t(x,0)=\psi(x), & x\in\mathbb{R},
\end{cases}
\end{align*}
其中 $a>0$ 是常数。
\end{example}

\begin{example}
利用 Fourier 变换求解问题
\begin{align*}
-\Delta u+u=f,\quad x\in\mathbb{R}^n.
\end{align*}
\end{example}

\begin{example}
利用 Fourier 变换求解 Schr\"{o}dinger 方程的初值问题
\begin{align*}
\begin{cases}
u_t-i\Delta u=0, & (x,t)\in\mathbb{R}^n\times\mathbb{R}_+,\\
u(x,0)=\varphi(x), & x\in\mathbb{R}^n.
\end{cases}
\end{align*}
\end{example}

\begin{example}
利用 Fourier 变换求解 Cahn-Hilliard 方程的初值问题
\begin{align*}
\begin{cases}
u_t+\Delta^2u=0, & (x,t)\in\mathbb{R}^n\times\mathbb{R}_+,\\
u(x,0)=\varphi(x), & x\in\mathbb{R}^n.
\end{cases}
\end{align*}
\end{example}

\begin{example}
假使 $u_1(s,t),u_2(s,t),\cdots,u_n(s,t)$ 满足热方程 $u_t-a^2u_{ss}=0$，则
\begin{align*}
u(x,t)=u(x_1,x_2,\cdots,x_n,t)=\prod_{k=1}^n u_k(x_k,t)
\end{align*}
满足热方程
\begin{align*}
u_t-a^2\Delta u=0,
\end{align*}
其中 $a$ 为正常数，$\Delta u=u_{x_1x_1}+u_{x_2x_2}+\cdots+u_{x_nx_n}$。
\end{example}

\begin{example}
假使 $u(x,t)\in C^3(\mathbb{R}_+^2)$ 满足热方程 $u_t-a^2u_{xx}=0$ $(a>0)$，则
\begin{enumerate}[(1)]
\item 对于任意 $\lambda\in\mathbb{R}$，$u_\lambda(x,t)=u(\lambda x,\lambda^2t)$ 满足热方程；
\item 函数 $v(x,t)=xu_x+2tu_t(x,t)$ 也满足热方程。
\end{enumerate}
\end{example}

\begin{example}
求解拟线性抛物方程的初值问题
\begin{align*}
\begin{cases}
u_t-a^2\Delta u+b|Du|^2=0, & (x,t)\in\mathbb{R}_+^{n+1},\\
u(x,0)=\varphi(x), & x\in\mathbb{R}^n,
\end{cases}
\end{align*}
其中 $a$ 为正常数，$b$ 为常数。（提示：作 Hopf-Cole 变换 $w=e^{-\frac{bu}{a^2}}$）
\end{example}

\begin{example}
求解粘性 Burgers 方程的初值问题
\begin{align*}
\begin{cases}
u_t-a^2u_{xx}+uu_x=0, & (x,t)\in\mathbb{R}_+^2,\\
u(x,0)=\varphi(x), & x\in\mathbb{R},
\end{cases}
\end{align*}
其中 $a$ 为常数。（提示：作函数变换 $v(x,t)=\int_{-\infty}^x u(y,t)\,\mathrm{d}y$，将问题化为上一问题）
\end{example}

\begin{example}
求解拟线性抛物方程的初值问题
\begin{align*}
\begin{cases}
u_t-\Delta u-\frac{U''(u)}{U'(u)}|Du|^2=0, & (x,t)\in\mathbb{R}_+^{n+1},\\
u(x,0)=\varphi(x), & x\in\mathbb{R}^n,
\end{cases}
\end{align*}
其中 $U:\mathbb{R}\to\mathbb{R}$ 是一个光滑函数且 $U'>0$。（提示：作函数变换 $v(x,t)=U[u(x,t)]$）
\end{example}

\begin{example}
求解拟线性抛物方程的初值问题
\begin{align*}
\begin{cases}
u_t-\Delta u+f(u)|Du|^2=0, & (x,t)\in\mathbb{R}_+^{n+1},\\
u(x,0)=\varphi(x), & x\in\mathbb{R}^n,
\end{cases}
\end{align*}
其中 $f:\mathbb{R}\to\mathbb{R}$ 是一个光滑函数。
\end{example}

\begin{example}
利用 Poisson 公式 \eqref{eq:poisson-1d} 证明 Weierstrass 逼近定理，即多项式在 $C([0,1])$ 上是稠密的。（提示：对给定 $\varphi\in C([0,1])$，延拓 $\varphi$ 到 $\mathbb{R}$ 上，求解初值问题
\begin{align*}
\begin{cases}
u_t-u_{xx}=0, & (x,t)\in\mathbb{R}_+^2,\\
u(x,0)=\varphi(x), & x\in\mathbb{R}.
\end{cases}
\end{align*}
其求解步骤如下：(1) 证明当 $t\to 0$ 时，$u(x,t)$ 一致地收敛到 $\varphi(x)$；(2) 将 $u(x,t)$ 用基本解 $K(x,t)$ 表示出来；(3) 用多项式逼近基本解 $K(x,t)$）
\end{example}

\begin{example}
利用 Poisson 公式 \eqref{eq:poisson-1d} 推导半无界问题
\begin{align*}
\begin{cases}
u_t-u_{xx}=0, & x>0,\ t>0,\\
u(x,0)=0, & x>0,\\
u(0,t)=g(t), & t>0
\end{cases}
\end{align*}
的显式解
\begin{align*}
u(x,t)=\frac{x}{\sqrt{4\pi}}\int_0^t \frac{1}{(t-s)^{\frac{3}{2}}}e^{-\frac{x^2}{4(t-s)}}g(s)\,\mathrm{d}s,
\end{align*}
其中 $g(0)=0$。（提示：令 $v(x,t)=u(x,t)-g(t)$，然后把 $v(x,t)$ 奇延拓到 $x<0$，再利用 Poisson 公式求解）
\end{example}

\begin{example}
设 $a$ 为正常数，$A_1,A_2$ 为常数。用分离变量法求解下列混合问题：
\begin{enumerate}[(1)]
\item $\begin{cases}
u_t-u_{xx}=u, & 0<x\leqslant\pi,\ t>0,\\
u(x,0)=\sin x, & 0\leqslant x\leqslant\pi,\\
u(0,t)=0,\quad u(\pi,t)=0, & t\geqslant 0;
\end{cases}$
\item $\begin{cases}
u_t-a^2u_{xx}=0, & 0<x\leqslant\pi,\ t>0,\\
u(x,0)=\cos x, & 0\leqslant x\leqslant\pi,\\
u_x(0,t)=0,\quad u_x(\pi,t)=0, & t\geqslant 0;
\end{cases}$
\item $\begin{cases}
u_t-a^2u_{xx}=0, & 0<x\leqslant l,\ t>0,\\
u(x,0)=x^2(l-x)^2, & 0\leqslant x\leqslant l,\\
u_x(0,t)=0,\quad u_x(l,t)=0, & t\geqslant 0;
\end{cases}$
\item $\begin{cases}
u_t-u_{xx}=0, & 0<x\leqslant\pi,\ t>0,\\
u(x,0)=0, & 0\leqslant x\leqslant\pi,\\
u_x(0,t)=A_1t,\quad u_x(\pi,t)=A_2t, & t\geqslant 0;
\end{cases}$
\item $\begin{cases}
u_t-u_{xx}=x\cos t, & 0<x\leqslant\pi,\ t>0,\\
u(x,0)=0, & 0\leqslant x\leqslant\pi,\\
u(0,t)=0,\quad u(\pi,t)=0, & t\geqslant 0.
\end{cases}$
\end{enumerate}
\end{example}

\begin{example}
设 $a>0$ 为常数。求下列混合问题的 Green 函数：
\begin{enumerate}[(1)]
\item $\begin{cases}
u_t-a^2u_{xx}=0, & 0<x<l,\ t>0,\\
u(0,t)=0,\quad u(l,t)=0, & t>0;
\end{cases}$
\item $\begin{cases}
u_t-a^2u_{xx}=0, & 0<x<l,\ t>0,\\
u_x(0,t)=0,\quad u_x(l,t)=0, & t>0;
\end{cases}$
\item $\begin{cases}
u_t-a^2u_{xx}=0, & 0<x<l,\ t>0,\\
u_x(0,t)=0,\quad u(l,t)=0, & t>0.
\end{cases}$
\end{enumerate}
\end{example}

\begin{example}
设 $a>0$ 为常数。利用上面混合问题的 Green 函数求解下列问题：
\begin{enumerate}[(1)]
\item $\begin{cases}
u_t-a^2u_{xx}=f(x,t), & 0<x<l,\ t>0,\\
u(x,0)=\varphi(x), & 0\leqslant x\leqslant l,\\
u(0,t)=u(l,t)=0, & t>0;
\end{cases}$
\item $\begin{cases}
u_t-a^2u_{xx}=f(x,t), & 0<x<l,\ t>0,\\
u(x,0)=\varphi(x), & 0\leqslant x\leqslant l,\\
u_x(0,t)=u_x(l,t)=0, & t>0;
\end{cases}$
\item $\begin{cases}
u_t-a^2u_{xx}=f(x,t), & 0<x<l,\ t>0,\\
u(x,0)=\varphi(x), & 0\leqslant x\leqslant l,\\
u_x(0,t)=0,\quad u(l,t)=0, & t>0.
\end{cases}$
\end{enumerate}
\end{example}

\begin{example}
设有长度为 $l$，两端绝热的均匀细杆，其初始温度为 $\varphi(x)$。考虑下列两种情形：
\begin{enumerate}[(1)]
\item 侧表面绝热；
\item 侧表面与周围介质发生热交换，且介质温度为常温 $u_0$。
\end{enumerate}
分别在上述两种情形下求杆上的温度分布 $u(x,t)$，并求 $t\to+\infty$ 时杆上的温度的极限分布。
\end{example}

\begin{example}
设有两段均匀且截面相同的细杆，两段杆的特性常数分别为 $c_1,\rho_1,k_1$ 和 $c_2,\rho_2,k_2$，长度分别为 $l_1$ 和 $l_2$。将它们连接成一组合杆。设组合杆的初始温度为 $\varphi(x)$，组合杆两端保持零度，侧表面绝热。试求组合杆上的温度分布。（提示：设两段杆的温度分布分别为 $u_1$ 和 $u_2$，则在连接处满足条件 $u_1=u_2$，$k_1\frac{\partial u_1}{\partial x}=k_2\frac{\partial u_2}{\partial x}$）
\end{example}

\begin{example}
证明附注 \eqref{eq:growth}。
\end{example}

\begin{example}
证明定理 \eqref{eq:green-sol-heat}。
\end{example}

\begin{example}
若 $v\in C^{2,1}(\Omega_T)$ 满足
\begin{align*}
v_t-a^2\Delta v\leqslant 0,\quad (x,t)\in\Omega_T,
\end{align*}
称 $v$ 在 $\Omega_T$ 上是热方程的下解，其中 $\Omega_T=\Omega\times(0,T]$，$a>0$ 为常数。
\begin{enumerate}[(1)]
\item 证明：
\begin{align*}
\max_{\overline{\Omega}_T}v(x,t)=\max_{\partial_p\Omega_T}v(x,t).
\end{align*}
\item 设 $\phi:\mathbb{R}\to\mathbb{R}$ 是光滑凸函数且 $u$ 在 $\Omega_T$ 上满足热方程。证明：$v=\phi(u)$ 在 $\Omega_T$ 上是热方程的下解。
\item 设 $u$ 在 $\Omega_T$ 上满足热方程。证明：$v=a^2|Du|^2+u_t^2$ 在 $\Omega_T$ 上是热方程的下解。
\end{enumerate}
\end{example}

在以下各题中，假设区域 $Q_T=\{(x,t)\mid 0<x<l,\ 0<t\leqslant T\}$，$\Gamma=\partial_p Q_T$ 是抛物边界。

\begin{example}
假设 $u\in C^{2,1}(Q_T)\cap C(\overline{Q}_T)$ 是热方程
\begin{align*}
u_t-u_{xx}=u,\quad (x,t)\in Q_T
\end{align*}
的非负解。假设存在正数 $M>0$，使得
\begin{align*}
u|_\Gamma\leqslant M.
\end{align*}
证明：
\begin{align*}
u(x,t)\leqslant Me^t,\quad (x,t)\in\overline{Q}_T.
\end{align*}
（提示：考虑 $w=e^{-t}u$ 满足的定解问题）
\end{example}

\begin{example}
设 $\mathcal{L}u=u_t-u_{xx}+|u_x|$。证明对于算子 $\mathcal{L}$，比较原理成立，即设 $u,v\in C^{2,1}(Q_T)\cap C(\overline{Q}_T)$，当 $\mathcal{L}u\leqslant \mathcal{L}v$，$u|_\Gamma\leqslant v|_\Gamma$ 时，则在 $\overline{Q}_T$ 上 $u(x,t)\leqslant v(x,t)$。
\end{example}

\begin{example}
设 $\mathcal{L}u=u_t-u_{xx}+u^3$。证明对于算子 $\mathcal{L}$，比较原理成立。

（提示：考虑 $w(x,t)=u(x,t)-v(x,t)$ 满足的定解问题）
\end{example}

\begin{example}
假设 $u\in C^{2,1}(\overline{Q}_T)$，$u_t\in C^{2,1}(Q_T)$ 且满足定解问题
\begin{align*}
\begin{cases}
u_t-u_{xx}=f(x,t), & (x,t)\in Q_T,\\
u(x,0)=\varphi(x), & 0\leqslant x\leqslant l,\\
u(0,t)=u(l,t)=0, & 0\leqslant t\leqslant T,
\end{cases}
\end{align*}
则
\begin{align*}
\max_{\overline{Q}_T}|u_t(x,t)|\leqslant C(\|f\|_{C^1(Q_T)}+\|\varphi''\|_{C[0,l]}),
\end{align*}
其中 $C$ 仅依赖于 $T$。
\end{example}

\begin{example}
假设 $u\in C^{1,0}(\overline{Q}_T)\cap C^{2,1}(Q_T)$ 且满足定解问题
\begin{align*}
\begin{cases}
u_t-u_{xx}=0, & (x,t)\in Q_T,\\
u(x,0)=\varphi(x), & 0\leqslant x\leqslant l,\\
u(0,t)=u(l,t)=0, & 0\leqslant t\leqslant T,
\end{cases}
\end{align*}
则
\begin{enumerate}[(1)]
\item $\max_{(0,T)}|u_x(0,t)|\leqslant C$，$\max_{(0,T)}|u_x(l,t)|\leqslant C$，其中 $C$ 仅依赖于 $\|\varphi\|_{C^1[0,l]}$；
\item 设 $u_x\in C^{2,1}(Q_T)$，则 $\max_{\overline{Q}_T}|u_x(x,t)|\leqslant \tilde{C}$，其中 $\tilde{C}$ 也仅依赖于 $\|\varphi\|_{C^1[0,l]}$。
\end{enumerate}
\end{example}

\begin{example}
假设 $u,u_x\in C(\overline{Q}_T)\cap C^{2,1}(Q_T)$ 且满足定解问题
\begin{align*}
\begin{cases}
u_t-u_{xx}=f(x,t), & (x,t)\in Q_T,\\
u(x,0)=\varphi(x), & 0\leqslant x\leqslant l,\\
[-u_x+\alpha u]|_{x=0}=g_1(t), & 0\leqslant t\leqslant T,\\
[u_x+\beta u]|_{x=l}=g_2(t), & 0\leqslant t\leqslant T,
\end{cases}
\end{align*}
其中 $\alpha\geqslant 0,\ \beta\geqslant 0$。试给出 $\max_{\overline{Q}_T}|u_x|$ 的估计。
\end{example}

\begin{example}
记 $Q_T^l=\{0<x<l,\ 0<t\leqslant T\}$，并设 $u^l\in C(\overline{Q}_T^l)\cap C^{2,1}(Q_T^l)$ 且满足定解问题
\begin{align*}
\begin{cases}
u_t^l-u_{xx}^l=0, & (x,t)\in Q_T^l,\\
u^l(x,0)=0, & 0\leqslant x\leqslant l,\\
u^l(0,t)=g(t),\quad u^l(l,t)=0, & 0\leqslant t\leqslant T,
\end{cases}
\end{align*}
其中 $g(t)\geqslant 0$。证明：$u^l(x,t)$ 关于 $l$ 是递增的，即对于 $l_1<l_2$，
\begin{align*}
u^{l_1}(x,t)\leqslant u^{l_2}(x,t),\quad (x,t)\in Q_T^{l_1}.
\end{align*}
\end{example}

\begin{example}
设 $u\in C^{1,0}(\overline{Q}_T)\cap C^{2,1}(Q_T)$ 且满足定解问题
\begin{align*}
\begin{cases}
u_t-u_{xx}=0, & (x,t)\in Q_T,\\
u(x,0)=0, & 0\leqslant x\leqslant l,\\
[u_x+h(u_0-u)]|_{x=0}=0,\quad u|_{x=l}=0, & 0\leqslant t\leqslant T,
\end{cases}
\end{align*}
其中 $h,u_0$ 为正常数。证明：
\begin{enumerate}[(1)]
\item $0\leqslant u(x,t)\leqslant u_0,\ (x,t)\in\overline{Q}_T$；
\item $u=u_h(x,t)$ 关于 $h$ 单调递增。
\end{enumerate}
\end{example}

\begin{example}
假设 $u\in C(\overline{Q}_T)\cap C^{2,1}(Q_T)$ 且满足定解问题
\begin{align*}
\begin{cases}
u_t-u_{xx}=-u^2+bu, & (x,t)\in Q_T,\\
u(x,0)=\varphi(x), & 0\leqslant x\leqslant l,\\
u(0,t)=u(l,t)=0, & 0\leqslant t\leqslant T,
\end{cases}
\end{align*}
其中 $b=b(x,t)\in C(\overline{Q}_T)$，$\varphi\in C[0,l]$，$\varphi(x)\geqslant 0$。证明：
\begin{align*}
0\leqslant u(x,t)\leqslant M\max_{0\leqslant x\leqslant l}\varphi(x),
\end{align*}
其中 $M$ 只依赖于 $T$，$\max_{\overline{Q}_T}|b(x,t)|$。
\end{example}

\begin{example}
证明：初值问题
\begin{align*}
\begin{cases}
u_t-a(x,t)u_{xx}+b(x,t)u_x+c(x,t)u=f(x,t), & (x,t)\in\mathbb{R}_+^2,\\
u(x,0)=\varphi(x), & x\in\mathbb{R}
\end{cases}
\end{align*}
的有界解是惟一的，其中 $a(x,t)\geqslant a_0>0$，$c(x,t)\geqslant 0$，且 $a(x,t)$，$b(x,t)$ 和 $c(x,t)$ 有界。（提示：首先在区域 $Q_T^L=(-L,L)\times(0,T]$ 上证明比较原理成立，然后考虑辅助函数 $w=\frac{M}{L^2}e^t(x^2+2At+B^2)$，其中 $|u(x,t)|\leqslant M$，$|a(x,t)|\leqslant A$，$|b(x,t)|\leqslant B$）
\end{example}

\begin{example}
假设 $u\in C^{2,1}(\overline{Q}_T)$ 且满足定解问题
\begin{align*}
\begin{cases}
u_t-u_{xx}=f(x,t), & (x,t)\in Q_T,\\
u(x,0)=\varphi(x), & 0\leqslant x\leqslant l,\\
u(0,t)=u(l,t)=0, & 0\leqslant t\leqslant T.
\end{cases}
\end{align*}
证明：
\begin{align*}
\sup_{0\leqslant t\leqslant T}\int_0^l u_x^2(x,t)\,\mathrm{d}x+\int_0^T\int_0^l u_t^2(x,t)\,\mathrm{d}x\,\mathrm{d}t
\leqslant 2\left\{\int_0^l[\varphi'(x)]^2\,\mathrm{d}x+\int_0^T\int_0^l f^2(x,t)\,\mathrm{d}x\,\mathrm{d}t\right\}.
\end{align*}
\end{example}

\begin{example}
假设 $u\in C^{1,0}(\overline{Q}_T)\cap C^{2,1}(Q_T)$ 且满足定解问题
\begin{align*}
\begin{cases}
u_t-a^2u_{xx}=f(x,t), & (x,t)\in Q_T,\\
u(x,0)=\varphi(x), & 0\leqslant x\leqslant l,\\
[-u_x+\alpha u]|_{x=0}=[u_x+\beta u]|_{x=l}=0, & 0\leqslant t\leqslant T,
\end{cases}
\end{align*}
其中 $a>0,\ \alpha\geqslant 0,\ \beta\geqslant 0$ 为常数。证明：
\begin{align*}
\sup_{0\leqslant t\leqslant T}\int_0^l u^2(x,t)\,\mathrm{d}x+\int_0^T\int_0^l u_x^2(x,t)\,\mathrm{d}x\,\mathrm{d}t
\leqslant M\left[\int_0^l\varphi^2(x)\,\mathrm{d}x+\int_0^T\int_0^l f^2(x,t)\,\mathrm{d}x\,\mathrm{d}t\right],
\end{align*}
其中 $M$ 只依赖于 $T$ 和 $a$。
\end{example}

\begin{example}
假设 $u\in C^{1,0}(\overline{Q}_T)\cap C^{2,1}(Q_T)$ 且满足定解问题
\begin{align*}
\begin{cases}
u_t-a^2u_{xx}+b(x,t)u_x+c(x,t)u=f(x,t), & (x,t)\in Q_T,\\
u(x,0)=\varphi(x), & 0\leqslant x\leqslant l,\\
u(0,t)=u(l,t)=0, & 0\leqslant t\leqslant T,
\end{cases}
\end{align*}
其中 $a$ 为正常数。证明：
\begin{align*}
\sup_{0\leqslant t\leqslant T}\int_0^l u^2(x,t)\,\mathrm{d}x+\int_0^T\int_0^l u_x^2(x,t)\,\mathrm{d}x\,\mathrm{d}t
\leqslant M\left[\int_0^l\varphi^2(x)\,\mathrm{d}x+\int_0^T\int_0^l f^2(x,t)\,\mathrm{d}x\,\mathrm{d}t\right],
\end{align*}
其中 $M$ 只依赖于 $T$，$a$，$B$ 和 $C$，$B=\sup\{|b(x,t)|\}$，$C=\sup\{|c(x,t)|\}$ 是非负常数。
\end{example}



\end{document}